Für die ganze Höhe gilt $h = \frac{1}{2}g\,T^2$, für die halbe Höhe $\frac{h}{2} = \frac{1}{2}g\,t^2$. Teilt man die zweite Gleichung durch die erste, fällt $g$ heraus:
\begin{align}
 \frac{1}{2} &= \frac{t^2}{T^2} \quad\Rightarrow\quad t = \thalF \\
   &= \frac{\Tg}{\sqrt{2}} = \thal \approx \result{\thalP{0}{2}}
\end{align}
Das ist der Bruchteil $\frac{1}{\sqrt{2}} \approx \result{\frP{0}{2}}$ der Fallzeit, unabhängig von der Höhe des Turms. Der Stein braucht für die obere Hälfte länger als für die untere, weil er oben noch langsam ist.
